\section{A Galilean Jewish fisherman}\label{sec:peter-world}

Peter belonged to the Jewish world of Roman Palestine. Galilee in his youth
was governed by Herod Antipas under Roman supremacy; Judea passed from Herod's
dynasty to direct Roman prefects. Taxation, patronage, local elites, imperial
power, and the presence of armed coercion were part of ordinary life. So were
synagogue, Sabbath, pilgrimage, purity, Scripture, prayer, and hope in the God
of Israel. To call Peter a ``Christian'' is unavoidable shorthand for his
later vocation, but it must not conceal that he followed Israel's Messiah as
a Jew and argued about the admission of Gentiles within a movement born
inside Israel.

\subsection{Bethsaida, Capernaum, and the fishing household}

John 1:44 calls Peter and his brother Andrew natives of Bethsaida.
Mark 1:21--31 places a house shared by Simon and Andrew in Capernaum and reports that
Simon's mother-in-law lay ill there. The accounts are compatible with an
origin in one lakeside settlement and a later household in another, but the
sources do not explain the move. Luke 5 portrays Simon working with partners,
including James and John; Mark 1 shows hired workers in Zebedee's boat. This
was labor embedded in households, partnerships, markets, tolls, and the
political economy of the lake. ``Fisherman'' should evoke arduous skilled
work, not either romantic destitution or unproved commercial wealth.

The mother-in-law establishes that Peter was married. Paul later asks whether
he and Barnabas lack the right to travel with a believing wife, ``as do the
other apostles and the brothers of the Lord and Cephas'' (1 Cor 9:5). The
verse is unusually valuable because it is early, incidental, and directed to
a practical dispute rather than composed to supply biography. Whether the
woman accompanied Peter on the particular mission Paul knew, and whether she
was the same wife as in Galilee, are reasonable but not demonstrable
inferences. Clement of Alexandria much later told of Peter encouraging his
wife at martyrdom; the story may preserve a remembered married apostolate, but
it does not disclose her name or establish the scene as contemporary report.

\subsection{Simon son of John---or Jonah}

The sources vary in the father's name. John calls him Simon son of John; the
Matthean form in 16:17 is traditionally rendered Simon Bar-Jona. This may
reflect related Semitic forms or a textual and translational history, but it
does not authorize a confident family genealogy. Andrew is consistently his
brother. No canonical source names Peter's mother, children, or later
descendants. Petronilla, venerated as a Roman saint and later called his
daughter in legend, cannot be inserted into the first-century household on
that basis.

Peter's languages and literacy likewise require restraint. A Galilean engaged
in fishing and travel could have encountered Aramaic, Hebrew in religious
settings, and some Greek. Acts 4:13 calls Peter and John \emph{agrammatoi kai
idiotai}, usually rendered uneducated and ordinary; the expression denies
elite scribal formation, not intelligence or all literacy. The polished Greek
of First Peter is consequently relevant to authorship, but it cannot by itself
prove that Peter never learned, dictated, or commissioned Greek. First Peter
5:12 explicitly mentions writing ``through Silvanus,'' a phrase whose precise
secretarial force is debated.

\subsection{The call in four canonical portraits}

Mark and Matthew place the summons beside the Sea of Galilee: Jesus calls
Simon and Andrew from their nets to become fishers of people. Luke expands
the scene into a miraculous catch, Simon's confession of sinfulness, and a
commission not to fear. John instead introduces Andrew as a disciple who
brings Simon to Jesus; Jesus immediately names him Cephas. These are not four
camera angles that must be spliced into a minute-by-minute sequence. Mark
emphasizes the authority of the summons and immediate response. Luke frames
mission through abundance, fear, and mercy. John places the naming at the
beginning of a testimony-driven gathering of disciples.

Their common center is sturdy: Simon belonged to the first circle of Jesus'
Galilean followers; Andrew was associated with his entrance into that circle;
his ordinary work became a governing metaphor for mission; and the name Rock
expressed a vocation given by Jesus, not a later self-promotion. The very
differences preserve how early communities contemplated that vocation.

\evidenceaxis{Peter was a married Jewish fisherman connected with Bethsaida
and a household in Capernaum.}{Mark 1:16--31; Luke 5:1--11; John 1:40--44; 1
Cor 9:5.}{The sources do not yield his birth date, wealth, formal education,
wife's name, or a continuous domestic biography.}
